% Bundle written by loom. Statements of the closure precede the key; ids are shown as comments before each environment.
\documentclass{amsart}
\usepackage{amsmath,amssymb,amsthm}
\usepackage{loom}

\theoremstyle{plain}
\newtheorem{theorem}{Theorem}[section]
\newtheorem{lemma}[theorem]{Lemma}
\newtheorem{proposition}[theorem]{Proposition}
\theoremstyle{definition}
\newtheorem{definition}[theorem]{Definition}
\theoremstyle{remark}
\newtheorem{remark}[theorem]{Remark}

\newcommand{\Fix}{\operatorname{Fix}}

\title{Widgets and their fixed loci}
\author{The loom demo}
\begin{document}
\section*{Bundle for dm-0003}

% id: dm-0003
\begin{theorem}[Main]\label{dm-0003}
Let $(X,\sigma)$ be a widget. Then $|\Fix(\sigma)| \equiv |X| \pmod 2$, so $\Fix(\sigma)$ is nonempty when $|X|$ is odd; and $\Fix(\sigma)$ is closed in every Hausdorff topology on $X$.
\end{theorem}

% proof: dm-0003/proof
\begin{proof}
\uses{dm-0002}
By Lemma~\ref{lem:orbits} the set $X$ is a disjoint union of orbits of size one or two, so $|X| \equiv |\Fix(\sigma)| \pmod 2$, which gives the parity statement. For closedness, $\Fix(\sigma)$ is the preimage of the diagonal under $(\mathrm{id},\sigma)\colon X\to X\times X$, and the diagonal is closed when $X$ is Hausdorff; the same argument in the algebraic setting is \cite[Proposition 3.2]{Man12}.
\end{proof}

\end{document}
